\documentclass[12pt,a4paper]{article}
        \makeatletter
        \def\input@path{{../}}
        \makeatother
        \usepackage{almeria-fichas}

        \FichaMeta{Sesion 20}{Exposicion y cierre}{Bloque 4 · Producto final}{Presentar la guia final, valorar otras exposiciones y cerrar el proyecto con una reflexion sobre lo aprendido.}{Procesos}{Recordar, Evaluar, Crear}

        \begin{document}

        \FichaCabecera

        \begin{materialesbox}
        \begin{itemize}
    \item Ficha impresa
    \item Lapiz
    \item Guia final del equipo
    \item Rubrica de observacion de exposiciones
\end{itemize}
        \end{materialesbox}

        \begin{pasosbox}
        \begin{enumerate}
    \item Escucha las exposiciones con respeto.
    \item Toma notas breves mientras observas.
    \item Usa la rubrica para valorar con justicia.
    \item Cierra la ficha con una reflexion personal.
\end{enumerate}
        \end{pasosbox}

        \begin{recordatoriobox}
        \begin{itemize}
    \item Exponer no es solo leer: hay que explicar.
    \item Una evaluacion justa usa evidencias observables.
    \item La reflexion final ayuda a consolidar el aprendizaje.
\end{itemize}
        \end{recordatoriobox}

        \begin{apoyobox}
        \begin{itemize}
    \item Anota palabras clave mientras escuchas.
    \item Si valoras una exposicion, justifica con un ejemplo concreto.
    \item En la reflexion final, usa frases cortas y sinceras.
\end{itemize}
        \end{apoyobox}

        \BloomSection{recordar}

\BloomExercise{recordar}{1}{Vocabulario esencial}{Escribe \textbf{cuatro palabras clave} relacionadas con \emph{Exposicion y cierre}. Despues escribe al lado una idea breve sobre cada una.}{6}

\BloomExercise{recordar}{2}{Datos o elementos basicos}{Anota \textbf{tres elementos importantes} de esta sesion. Ten en cuenta que la dimension procesos prioriza tomar decisiones, organizar el trabajo y reflexionar sobre como se actua.}{6}

\BloomSection{evaluar}

\BloomExercise{evaluar}{1}{Observacion de exposiciones}{Evalua tres exposiciones de la clase con dos criterios: claridad matematica y calidad visual. Justifica cada valoracion.}{9}

\BloomExercise{evaluar}{2}{Que diferencia a un buen trabajo}{Escribe que diferencia, en tu opinion, una guia notable de una sobresaliente.}{7}

\BloomSection{crear}

\BloomExercise{crear}{1}{Nueva ciudad}{Imagina que ahora debes empezar una guia matematica sobre otra ciudad. Crea la propuesta de sus dos primeras paginas.}{10}

\BloomExercise{crear}{2}{Cierre personal}{Escribe un texto final breve: que he aprendido, que me ha costado mas y donde he visto matematicas en la ciudad.}{8}

        \begin{evidenciabox}
        \begin{itemize}
    \item Valorar varias exposiciones con criterio.
    \item Reconocer rasgos de una buena guia matematica.
    \item Reflexionar sobre lo aprendido durante el proyecto.
\end{itemize}
        \end{evidenciabox}

        \end{document}
